\chapter{Continuous Functions}
\label{chap:continuous_functions}

Continuity is the central structural morphism between metric spaces. This chapter provides foundational characterisations of continuous maps, uniform continuity, homeomorphisms, and extension theorems.

\section{Definition and Characterisations of Continuity}
\label{sec:continuity_def}

\begin{definition}[Continuous Function at a Point]
Let $(X, d_X)$ and $(Y, d_Y)$ be metric spaces. A function $f: X \to Y$ is called \textbf{continuous at a point} $x_0 \in X$ if for every $\eps > 0$, there exists $\delta > 0$ such that:
\[
d_X(x, x_0) < \delta \implies d_Y(f(x), f(x_0)) < \eps.
\]
In ball notation, $f(B_X(x_0, \delta)) \subseteq B_Y(f(x_0), \eps)$. A function $f$ is \textbf{continuous on $X$} if it is continuous at every point $x \in X$.
\end{definition}

\begin{theorem}[Topological and Sequential Characterisations of Continuity]
Let $f: (X, d_X) \to (Y, d_Y)$ be a mapping. The following statements are equivalent:
\begin{enumerate}[label=\rm{(\roman*)}]
    \item $f$ is continuous on $X$.
    \item For every open set $V \subseteq Y$, the inverse image $f^{-1}(V)$ is an open set in $X$.
    \item For every closed set $F \subseteq Y$, the inverse image $f^{-1}(F)$ is a closed set in $X$.
    \item For every sequence $(x_n)$ in $X$ converging to $x$, the sequence $(f(x_n))$ converges to $f(x)$ in $Y$.
    \item For every subset $A \subseteq X$, $f(\overline{A}) \subseteq \overline{f(A)}$.
\end{enumerate}
\end{theorem}

% TikZ Figure for Inverse Image Characterisation of Continuous Mapping
\begin{figure}[htbp]
\centering
\begin{tikzpicture}[scale=1.3]
    % Domain X
    \draw[thick, fill=mainblue!10, draw=mainblue, smooth cycle] 
        plot coordinates {(0,0) (2.5,0.8) (3,-1.5) (1,-2.5) (-0.5,-1)};
    \node[mainblue] at (0.2, 0.5) {$(X, d_X)$};
    
    % Preimage f^-1(V)
    \draw[thick, dashed, fill=darkred!20, draw=darkred, smooth cycle] 
        plot coordinates {(1,-0.5) (2,-0.2) (2.2,-1.5) (1.2,-1.8)};
    \node[darkred] at (1.6, -1) {$f^{-1}(V)$};
    
    % Codomain Y
    \draw[thick, fill=darkgreen!10, draw=darkgreen, smooth cycle] 
        plot coordinates {(6,0) (8.5,0.8) (9,-1.5) (7,-2.5) (5.5,-1)};
    \node[darkgreen] at (6.2, 0.5) {$(Y, d_Y)$};
    
    % Open Set V
    \draw[thick, fill=darkred!30, draw=darkred, smooth cycle] 
        plot coordinates {(7,-0.4) (8,-0.2) (8.2,-1.4) (7.1,-1.7)};
    \node[darkred] at (7.6, -0.9) {$V$};
    
    % Arrow f
    \draw[->, >=stealth, ultra thick, mainblue] (3.2, -0.5) -- node[above] {$f$} (5.3, -0.5);
    \draw[->, >=stealth, thick, dashed, darkred] (5.3, -1.2) -- node[below] {$f^{-1}$} (3.2, -1.2);
\end{tikzpicture}
\caption{Continuous mapping $f: X \to Y$: the inverse image $f^{-1}(V)$ of any open set $V \subseteq Y$ is an open set in $X$.}
\label{fig:continuous_mapping}
\end{figure}


\section{Uniform Continuity and Lipschitz Continuity}
\label{sec:uniform_continuity}

\begin{definition}[Uniform Continuity]
A function $f: (X, d_X) \to (Y, d_Y)$ is \textbf{uniformly continuous} on $X$ if for every $\eps > 0$, there exists $\delta > 0$ (depending only on $\eps$, independent of $x$) such that:
\[
d_X(x_1, x_2) < \delta \implies d_Y(f(x_1), f(x_2)) < \eps \quad \forall x_1, x_2 \in X.
\]
\end{definition}

\begin{definition}[Lipschitz Continuity]
A function $f: (X, d_X) \to (Y, d_Y)$ is called \textbf{Lipschitz continuous} if there exists a constant $K \ge 0$ such that:
\[
d_Y(f(x_1), f(x_2)) \le K \, d_X(x_1, x_2) \quad \forall x_1, x_2 \in X.
\]
If $K < 1$, $f$ is called a \textbf{contraction mapping}.
\end{definition}


\section{Homeomorphisms and Equivalent Metrics}
\label{sec:homeomorphism}

\begin{definition}[Homeomorphism]
A function $f: (X, d_X) \to (Y, d_Y)$ is called a \textbf{homeomorphism} if:
\begin{enumerate}[label=\rm{(\roman*)}]
    \item $f$ is bijective (one-to-one and onto).
    \item $f$ is continuous.
    \item $f^{-1}$ is continuous.
\end{enumerate}
If a homeomorphism exists between $X$ and $Y$, the spaces are said to be \textbf{homeomorphic} ($X \cong Y$).
\end{definition}

\begin{definition}[Uniformly Equivalent Metrics]
Two metrics $d_1$ and $d_2$ on the same set $X$ are called \textbf{equivalent} if they generate the same open sets, and \textbf{uniformly equivalent} if the identity map $i: (X, d_1) \to (X, d_2)$ is a uniform homeomorphism.
\end{definition}
